
%%%%%%%%%%%%%%%%%%%%%%%%%%%%%%%%%%%%%%%%%%%%%%%%%%%%%%%%%%%%%
\subsubsection{具体分析}

% ==== 第1段：问题类型 + 大致方法 ====
% 示例（A092风格）：问题一要求我们求解给定条件下XXX的XXX、XXX以及XXX。
问题一要求我们...，属于...类问题，本问采用...方法进行求解。

% ==== 第2段：完整步骤串联（首先...；接着...；然后...；最后...）====
首先，我们可以通过...来确定...；接着，通过...我们可以由...推得...，进而推得其...；然后，基于上述信息我们可以计算...，从而得出...；最后，结合...以及...的计算公式，代入...等数据，我们就可以求出...，进而求出...。

% ==== 第3段：引出流程图 ====
具体分析流程如图\ref{fig:analysis1_flow}所示。

% ==== 流程图（只改{}中文字，每框≤7字）====
\begin{figure}[H]
  \begin{center}
    \begin{tikzpicture}[x=1cm, y=0.7cm]
      \node[box] (n11) at (0, 0) {步骤1};
      \node[box] (n12) at (5, 0) {步骤2};
      \node[box] (n13) at (10, 0) {步骤3};
      \node[box] (n22) at (5, -3) {步骤5};
      \node[box] (n21) at (0, -3) {步骤4};
      \node[box] (n23) at (10, -3) {步骤6};
      \node[box] (n31) at (0, -6) {步骤7};
      \node[box] (n32) at (5, -6) {步骤8};
      \node[box] (n33) at (10, -6) {步骤9};
      \draw[arrow] (n11) -- (n12);
      \draw[arrow] (n12) -- (n13);
      \draw[arrow] (n23) -- (n22);
      \draw[arrow] (n22) -- (n21);
      \draw[arrow] (n13) -- ++(0,-1.0) -| (n23);
      \draw[arrow] (n21) -- (n31);
      \draw[arrow] (n31) -- (n32);
      \draw[arrow] (n32) -- (n33);
    \end{tikzpicture}
    \caption{问题一分析流程图}
    \label{fig:analysis1_flow}
  \end{center}
\end{figure}

% ============================================================
\subsubsection{模型准备}

% ==== 引文 ====
在进行建模之前，需要对数据进行预处理并选取合适的算法。本节将从数据预处理和算法选择两方面进行准备。

% ==== 数据预处理（仅问题1写此段）====
原始数据来源于...，包含...条记录和...个字段。经检查发现以下问题：...（缺失值/量纲不统一/异常值等）。

针对上述问题，按以下步骤进行预处理：
第一步，...。
第二步，...（如标准化，公式如下）：
\begin{equation}
x' = \frac{x - x_{\min}}{x_{\max} - x_{\min}}
\end{equation}
其中，$x_{\min}$、$x_{\max}$分别为最小值和最大值。

预处理后，数据规模为...行$\times$...列，各指标统计量如下表所示。

\begin{longtable}{>{\centering\arraybackslash}p{0.12\textwidth}>{\centering\arraybackslash}p{0.18\textwidth}>{\centering\arraybackslash}p{0.18\textwidth}>{\centering\arraybackslash}p{0.18\textwidth}>{\centering\arraybackslash}p{0.18\textwidth}}
  \caption{数据预处理前后统计量对比}
  \label{tab:预处理对比} \\
  \toprule
  指标 & 预处理前均值 & 预处理前标准差 & 预处理后均值 & 预处理后标准差 \\
  \midrule
  \endfirsthead
  \caption{数据预处理前后统计量对比（续）} \\
  \toprule
  指标 & 预处理前均值 & 预处理前标准差 & 预处理后均值 & 预处理后标准差 \\
  \midrule
  \endhead
  \bottomrule
  \endfoot
  ... & ... & ... & ... & ... \\
\end{longtable}

综上所述，数据预处理完成，为后续建模提供了清洁的输入数据。

% ==== 算法选择（A092风格：问题本质+候选算法+对比表+结论）====
该问题本质上是一个...问题，需要在...条件下...。常用求解算法包括...算法、...算法、...算法等。本节从...、...、...三方面进行评估。

\begin{longtable}{>{\centering\arraybackslash}p{0.12\textwidth}>{\centering\arraybackslash}p{0.18\textwidth}>{\centering\arraybackslash}p{0.18\textwidth}>{\centering\arraybackslash}p{0.18\textwidth}>{\centering\arraybackslash}p{0.18\textwidth}}
  \caption{算法能力对比}
  \label{tab:算法对比1} \\
  \toprule
  算法名称 & 全局搜索能力 & 收敛速度 & 约束处理能力 & 求解稳定性 \\
  \midrule
  \endfirsthead
  \caption{算法能力对比（续）} \\
  \toprule
  算法名称 & 全局搜索能力 & 收敛速度 & 约束处理能力 & 求解稳定性 \\
  \midrule
  \endhead
  \bottomrule
  \endfoot
  算法A（缩写） & 优 & 优 & 良 & 良 \\
  算法B（缩写） & 良 & 中 & 优 & 中 \\
  算法C（缩写） & 中 & 良 & 良 & 中 \\
\end{longtable}

通过对比可以看出，...算法在...方面表现最优，...，因此本文选择...算法作为核心求解方法。

综上所述，本节完成了数据预处理和算法选型，接下来将建立数学模型并进行求解。
